\subsubsection{Image Buffers}
\label{subsubsec:buffers}
% \todo[inline]{Citations}

Two different buffers with the same size (\texttt{BS}) are used.
A Baumer image buffer and an OpenCV cirucular buffer.

% For the above reasons, two image buffers are used.
% The received images are stored in Baumer \texttt{BGAPI2::Buffer} objects.
% To save the images later on, a circular buffer for OpenCV matrices (\texttt{cv::Mat}) is used.
% This is documented in section \ref{subsubsec:buffers}.
% The size of those image buffers (\texttt{BS}) determines the max. amount of valid frames $N_\text{max}$ that can be captured.

% Buffer handling by Baumer => large buffer sizes (here not a problem), could be handled by ourselves!
  % handle this by ourselves on the Ultra96 board, due to memory size constrains
\paragraph{Baumer Image Buffer}
The received images can either be stored in \texttt{BGAPI2::Buffer}
For convenience, the image buffer is hadeled by the Baumer GAPI.
% https://www.baumer.com/ch/en/product-overview/industrial-cameras-image-processing/software/baumer-gapi-sdk/c/14174

% The function BGAPI2::Buffer::Buffer() creates a buffer the size of the maximum
% payload of the camera that is connected. This enables switching between the pixel formats and various image details without the need for size checking or recreating the buffer,
% since all formats will fit within the buffer. 
% If a lot of buffer space is required for small image details (ROI); this creates a high memory
% requirement, although only a small portion of this is actually used.

% Alternatively the user may allocate the buffer memory on its own using the function BGAPI2::Buffer:: Buffer( void *pUserBuffer, bo_uint64
% uUserBufferSize, void *pUserObj);
% The responsibility to release the allocated memory then remains with the user.
% The buffer size can be adjusted to the required payload size. This enables many buffers
% with small image details (ROI).
% The payload size for the current camera parameter settings (ROI or pixel format) can
% be determined using the query pDevice->GetRemoteNode(„PayloadSize")-
% >GetInt().


\begin{table}[h]
  \caption{Baumer image buffer memory utilization of the pixel formats \texttt{BayerRG8} and \texttt{BGR8}}
  \label{tab:buffer_memory_utilization}
  \centering
  \begin{tabular}{lll}
    \toprule
     & \textbf{\texttt{BayerRG8}} & \textbf{\texttt{BGR8}} \\
    \midrule
    \textbf{Bytes per pixel} & \SI{1}{B} & \SI{3}{B} \\
    \textbf{Pixel count} & \SI{1310720}{px} & \SI{1310720}{px} \\
    \textbf{Pixel bytes} & \SI{1310720}{B} & \SI{3932160}{B} \\
    \textbf{Buffer size} & \SI{3932160}{B} & \SI{3932160}{B} \\
    \textbf{Utilization} & \SI{33}{\percent} & \SI{100}{\percent} \\
    \bottomrule
  \end{tabular}
\end{table}

\paragraph{OpenCV Circular Buffer}

% https://docs.opencv.org/4.1.1/d3/d63/classcv_1_1Mat.html#a51615ebf17a64c968df0bf49b4de6a3a

% OpenCV \texttt{Mat}

% opencv matrix circular buffer (same size as Baumer buffer but shallow [shallow-copy])
  % only pointer to the baumer pixel data
  % copying would require a lot of space and time (which is not much to play with), thus main reason time
% circular buffer graphic

% https://www.baumer.com/ch/en/service-support/know-how/technical-information-industrial-cameras/baumer-gapi-and-opencv/a/baumer-gapi-and-opencv

\begin{equation}
  I \equiv \text{FID} \pmod{\text{BS}}
  \label{eq:indexing}
\end{equation}

\begin{lstlisting}[style=C++]
  cv_buffer[frame_id % buff_size] = cv::Mat(height, width, CV_8UC1, (void *) pBufferFilled->GetMemPtr());
\end{lstlisting}

\begin{figure}[h]
  \centering
  \includegraphics[width=0.9\textwidth]{buffer}
  \caption{OpenCV matrices circular buffer}
  \label{fig:circular_buffer}
\end{figure}

% Since RAW formats often have to be converted into other target formats, such as BGR8 or
% BGR12, you can couple an external user buffer to this internal buffer, which then accepts
% the transformed data.

% Using pBuffer->GetMemPtr() accesses the RAW image and pBuffer
% ->GetUserObj() accesses the memory area in the user buffer. In the example, the
% external user buffer is of a size that equates to the target pixel format "BGR8", which
% requires 3 bytes per pixel.
